\documentclass[11pt]{article}
\usepackage[margin=0.78in]{geometry}
\usepackage{amsmath,amssymb,bm}
\usepackage{booktabs}
\usepackage{enumitem}
\usepackage{hyperref}
\hypersetup{colorlinks=true, linkcolor=black, urlcolor=black}
\setlength{\parindent}{0pt}
\setlength{\parskip}{0.55em}

\title{Cyrus Guided Notes: MIT 6.003 Signals and Systems}
\author{MIT course pack}
\date{2026-05-15}

\begin{document}
\maketitle

\section*{Course source}
\textbf{Official source:} \url{https://ocw.mit.edu/courses/6-003-signals-and-systems-fall-2011/}

\textbf{Why this course is in the system.} Prerequisite for signals, systems, convolution, frequency response, filtering, and state space.

\textbf{Learning output.} Write the time-domain to frequency-domain bridge in GoodNotes and add convolution visualization on the web.

\section*{Prerequisite checkpoint}
\begin{itemize}[leftmargin=1.4em]
  \item Complex numbers
  \item Calculus
  \item Differential equations
\end{itemize}

\section*{Formula checkpoint}
Do not memorize these first. Read each formula as a sentence: what changes, what is observed, what is optimized, and what output it produces.
\begin{align*}
  y(t)=x(t)*h(t)=\int_{-\infty}^{\infty}x(\tau)h(t-\tau)d\tau\\
  X(j\omega)=\int_{-\infty}^{\infty}x(t)e^{-j\omega t}dt\\
  H(s)=\frac{Y(s)}{X(s)}\\
  \dot{\bm{x}}=A\bm{x}+B\bm{u}
\end{align*}

\section*{Visual page}
Show an input signal sliding across an impulse response to create convolution.

\section*{GoodNotes pages}
\begin{itemize}[leftmargin=1.4em]
  \item Page MIT-S001: convolution
  \item Page MIT-S002: frequency response
\end{itemize}

\section*{Web interaction}
A slider for frequency that updates magnitude and phase.

\section*{Obsidian and Notion}
\begin{tabular}{p{0.22\linewidth}p{0.68\linewidth}}
\toprule
Layer & Output \\
\midrule
Obsidian & MIT EECS -> 6.003 signals systems \\
Notion & Signal exercise and formula status. \\
\bottomrule
\end{tabular}

\section*{First study move}
Open the official source, copy the prerequisite checkpoint into GoodNotes, then write one formula in your own words before watching or reading more.

\end{document}
